\chapter{Conclusion and Future Work}
\label{chap:conclusion}

\section{Conclusion}
\label{sec:conclusion}

This project designed, implemented, deployed, and empirically evaluated CORAL, a five-stage post-processing pipeline for low-resource Urdu ASR. The work was carried out across two iterations of the FAST-NUCES Final Year Project, and this final report consolidates both iterations into a single integrated narrative.

The defining technical contribution of CORAL is a \emph{split-merge-aware alignment algorithm} that treats word-boundary disagreement between ASR back-ends as a first-class phenomenon, classifying each aligned chunk as SAME, SPLIT, MERGE, or NOISE. Our empirical analysis on the 500-utterance Conversational Urdu sample shows that SPLIT and MERGE events together account for 36.5\% of all alignment events --- making segmentation disagreement the dominant inter-model divergence in Urdu ASR ensembles, ahead of pure lexical substitution. Existing fusion methods such as ROVER, designed for English-style space-delimited scripts, misinterpret these events as substitutions or insert--delete pairs, corrupting the voting signal. CORAL exposes the metadata and acts on it.

On the 2{,}995-utterance Common Voice Urdu evaluation, the deployed pipeline (\texttt{robust} category) reduces WER from 18.45\% to 14.34\% on the Seamless-Large source (a 22.3\% relative improvement), from 28.29\% to 19.97\% on Whisper-Large-v3 (29.4\% relative), from 40.44\% to 30.64\% on Whisper-Medium (24.2\% relative), and from 53.52\% to 39.67\% on Wav2Vec2-Urdu (25.9\% relative). On conversational Urdu, the full C7 pipeline (with the Stage~4 LLM refinement) reaches 10.6\% WER from a 19.8\% Whisper-Large-v3 baseline --- a \textbf{46.5\% relative improvement} that places CORAL slightly ahead of the closest concurrent system (Multi-ASR + SpeechLLM at 11.3\%) while avoiding the heavy cost of running a full SpeechLLM audio inference pass during alignment.

The seven-step ablation study on the Conversational Urdu benchmark demonstrates that every CORAL component makes a measurable, additive contribution. No single component alone explains the gain; the system is genuinely compositional. The Stage~4 LLM delivers the largest single-component gain (2.5 WER points), while Stage~0 Urdu normalisation --- a zero-cost, zero-risk fix --- contributes 1.9 WER points and confirms that a substantial fraction of apparent ASR errors are measurement artefacts caused by Arabic--Urdu Unicode-variant mismatches rather than genuine recognition failures.

The project also demonstrates the value of revisiting design decisions in light of empirical findings. The original FYP-1 architecture framed CORAL as a two-stage \emph{Generate-and-Refine} system in which an instruction-tuned LLM would arbitrate between confidence-annotated multi-model hypotheses. Iteration~1 evaluation revealed that confidence scores from heterogeneous architectures lived on incompatible scales, that synthetic hypotheses did not capture the dominant real-world failure modes (Unicode conflation and segmentation disagreement), and that the LLM was more useful as a refinement layer than as the primary fusion mechanism. The FYP-2 redesign therefore reframed the system as a five-stage post-processing pipeline (Stages~0--4) and dropped, replaced, or de-emphasised five FYP-1 elements documented in Section~\ref{ssec:dropped-components}: word-level confidence extraction from Whisper and Wav2Vec2, Expected Calibration Error as a primary metric, the synthetic 20-sample evaluation with calibration regimes, the Alif-1.0-8B-Instruct LLM, and the two-stage Generate-and-Refine framing itself. The FYP-2 system replaces these with deterministic algorithmic components (split-merge alignment, OOV detection, voting) followed by an LLM refinement layer that consumes structured metadata and is constrained by the upstream pipeline's authority limits.

Beyond the algorithmic contributions, the project delivers a complete deployable system: a FastAPI backend exposing three pipeline endpoints plus a model registry; a Next.js + React front-end implementing a four-pass user workflow with both file-upload and live-speech modes and animated alignment visualisations; Kaggle GPU notebooks that self-register through ngrok HTTPS tunnels and provide live ASR inference; a HuggingFace-hosted BK-tree and n-gram corpus loaded into DuckDB at server startup; and an open evaluation framework producing reproducible per-stage WER and CER metrics.

We answer the research hypothesis stated in Section~\ref{ssec:hypothesis} affirmatively: combining Urdu-specific Unicode normalisation, weighted split-merge-aware alignment of multiple ASR hypotheses, hybrid OOV detection and correction with BK-tree fuzzy lookup ranked by an n-gram language model, conservative position-wise voting, and targeted LLM refinement of the voted output \emph{is} sufficient to substantially reduce WER over the strongest single-model baseline on both read-speech and conversational Urdu, without requiring any fine-tuning of the underlying acoustic models.

\section{Future Work}
\label{sec:future-work}

The residual-error analysis (Section~\ref{ssec:residual-errors}) and the ablation study identify several concrete directions for follow-up research.

\paragraph{1. Confidence-weighted voting.}
All companion models currently receive equal voting weight. Model-specific WER priors used as confidence weights in Stage~3 voting --- in particular giving Whisper-Large or Seamless-Large a multiplicative weight $>$1 --- are estimated to yield a further $\sim$0.3 WER point reduction. This directly addresses the canonical ROVER failure mode observed on the read-speech benchmark, where three companions of WER 1.6$\times$--3.6$\times$ higher than the source collectively outvote the source on disagreement positions.

\paragraph{2. Code-switching-aware normalisation.}
Stage~0 currently strips all ASCII characters, deferring code-switched English to the LLM. A future mode will preserve English tokens and pass them through a bilingual alignment stage that explicitly handles Urdu--English code-switching, directly targeting the 21.5\% of residual errors attributable to code-switching identified in Table~\ref{tab:residual}.

\paragraph{3. Named-entity gazetteer.}
27.7\% of residual errors are proper nouns. A curated Urdu named-entity list (toponyms, person names, organisation names, brand names) integrated into the BK-tree lookup would directly address the dominant failure mode. The integration is straightforward: the gazetteer would be loaded as a high-priority secondary BK-tree consulted when the primary BK-tree does not return a high-confidence candidate.

\paragraph{4. LLM-guided top-$K$ OOV re-ranking.}
The current OOV correction is greedy top-1. Passing the top-3 BK-tree candidates to the Stage~4 LLM for context-aware re-ranking would reduce overcorrection errors (currently 7.2\% of residual errors) by giving the LLM the option to keep the original token when none of the candidates is a clear improvement.

\paragraph{5. Split/merge metadata utilisation in the voting stage.}
CORAL computes rich SAME/SPLIT/MERGE/NOISE metadata per token in Stage~1 but does not yet use it in the Stage~3 correction decision. Incorporating alignment type into the voting logic (e.g.\ different resolution strategies for SPLIT vs.\ SUBSTITUTION events) could reduce word-boundary errors further.

\paragraph{6. Word-level confidence from acoustic models.}
The FYP-1 architecture-specific confidence-extraction methodology was de-emphasised in FYP-2 because raw confidences from heterogeneous architectures lived on incompatible scales. A learned cross-architecture confidence calibration --- e.g.\ a TruCLeS-style~\cite{trucles2025} metric trained on a small labelled Urdu development set --- could provide reliable per-word uncertainty estimates that would improve both the voting stage and the LLM's arbitration.

\paragraph{7. Larger-scale conversational evaluation.}
The 500-clip Conversational Urdu sample used in the ablation study is a useful but limited test set. Scaling the evaluation to thousands of conversational clips --- particularly including more varied dialects, age groups, and acoustic conditions --- will firm up the empirical claims and may expose new failure modes.

\paragraph{8. Real-time streaming integration.}
CORAL is currently a batch-mode post-processor. A streaming variant that operates on partial hypotheses as they arrive from the ASR back-ends, emitting partial corrected transcripts with bounded latency, would unlock real-time captioning use cases.

\paragraph{9. Adaptation to other low-resource languages.}
The CORAL architecture is not Urdu-specific. The Stage~0 normalisation table, BK-tree corpus, and n-gram corpus would need to be replaced for each target language, but Stages~1--4 transfer directly. A natural next step is to validate the architecture on other low-resource languages of the South Asian subcontinent (Pashto, Sindhi, Punjabi).

\paragraph{10. Open release.}
The complete CORAL source code, evaluation TSV, BK-tree, and n-gram corpus are released under permissive licences in the project repository. We hope this enables further research on low-resource Urdu ASR post-processing.
